%%% =======================================================================
%%% ISLS Extension Phase 8 v1.0.0
%%% C27: The Foundry — Closed-Loop Software Fabrication
%%% Compile-Test-Fix Pipeline + Workspace Intelligence
%%% Author: Sebastian Klemm
%%% Date: 18 March 2026
%%% Status: Implementation Specification for Claude Code
%%% =======================================================================
\documentclass[11pt,a4paper,twoside]{article}

\usepackage[utf8]{inputenc}
\usepackage{amsmath,amssymb,amsthm}
\usepackage{algorithm,algorithmic}
\usepackage{booktabs,longtable,tabularx}
\usepackage{enumitem}
\usepackage{geometry}
\usepackage{hyperref}
\usepackage{cleveref}
\usepackage{xcolor}
\usepackage{fancyhdr}
\usepackage{listings}
\usepackage{titlesec}
\usepackage{tikz}
\usetikzlibrary{arrows.meta,positioning}

\geometry{margin=2.5cm,top=3cm,bottom=3cm,headheight=14pt}

\theoremstyle{definition}
\newtheorem{definition}{Definition}[section]
\newtheorem{axiom}{Axiom}[section]
\newtheorem{requirement}{Requirement}[section]
\newtheorem{invariant}{Invariant}[section]

\theoremstyle{plain}
\newtheorem{theorem}{Theorem}[section]

\theoremstyle{remark}
\newtheorem{remark}{Remark}[section]

\newcommand{\ISLS}{\textsf{ISLS}}
\newcommand{\FOUNDRY}{\textsf{FOUNDRY}}

\definecolor{sec-color}{HTML}{1e3a5f}

\titleformat{\section}
  {\Large\bfseries\color{sec-color}}
  {\thesection}{1em}{}

\lstset{
  basicstyle=\ttfamily\small,
  breaklines=true,
  frame=single,
  backgroundcolor=\color{gray!5},
  rulecolor=\color{gray!30},
}

\pagestyle{fancy}
\fancyhf{}
\fancyhead[LE,RO]{\ISLS{} Extension Phase 8 v1.0.0}
\fancyhead[RE,LO]{Sebastian Klemm}
\fancyfoot[C]{\thepage}

\hypersetup{
  colorlinks=true,
  linkcolor=blue!60!black,
  pdftitle={ISLS Phase 8 — The Foundry},
  pdfauthor={Sebastian Klemm},
}

\begin{document}

\begin{titlepage}
\centering
\vspace*{2.5cm}
{\Huge\bfseries \ISLS{} Extension Phase~8\\[0.3cm]
v1.0.0}\\[1.5cm]
{\Large C27: The \FOUNDRY{}}\\[1cm]
{\large Closed-Loop Software Fabrication\\[0.3cm]
Compile $\to$ Test $\to$ Fix $\to$ Until It Works\\[0.3cm]
+ Workspace Intelligence}\\[0.5cm]
{\normalsize The crate that makes the Forge\\
actually produce working software.}\\[2cm]
{\large Sebastian Klemm}\\[0.5cm]
{\large 18 March 2026}\\[1.5cm]

\begin{abstract}
\noindent
This document specifies \texttt{isls-foundry} (C27), the closed-loop
fabrication engine that bridges the gap between ``the Forge generates
code strings'' and ``the operator has a working, compilable, tested
software project on disk.''

The Foundry wraps the entire Forge pipeline (C21--C26) in a
build-test-fix loop: generated code is written to a real project
directory, compiled by the real toolchain (\texttt{cargo}), tested
by the real test runner, and --- if compilation or tests fail --- the
error output is fed back to the Oracle (C25) as a correction prompt.
The loop iterates until the project compiles cleanly, tests pass, or
the retry budget is exhausted.

Additionally, the Foundry provides Workspace Intelligence: the
ability to understand an existing codebase, identify what exists,
what is missing, and generate code that fits into the existing
project rather than creating a new one from scratch. This enables
incremental development: ``add an endpoint,'' ``implement this trait,''
``fix this module'' --- not just ``generate a whole new project.''

The Foundry is the layer that turns \ISLS{} from a formal system
that happens to generate code into a practical software development
tool.

\medskip\noindent\textbf{Status:} Implementation Specification.\\
\textbf{Parent:} ISLS v1.0.0 + Extensions Phase 1--7.\\
\textbf{Depends on:} C23 (Forge), C24 (Compose), C25 (Oracle),
  C26 (Templates).
\end{abstract}
\end{titlepage}

\tableofcontents
\newpage


%%% =====================================================================
\part{Why This Changes Everything}\label{part:why}

\section{The Gap}\label{sec:gap}

Without the Foundry, the Forge pipeline is:
\begin{lstlisting}
DecisionSpec -> PMHD -> ArtifactIR -> Oracle -> String -> Crystal
\end{lstlisting}

The output is a string that \emph{claims} to be Rust code. But nobody
checked. The string might have syntax errors, missing imports, wrong
types, incompatible trait bounds, or simply not compile. The 8-gate
cascade validates the crystal's \emph{formal} properties (hash
integrity, evidence chain, etc.), not the \emph{semantic} correctness
of the code inside it.

With the Foundry, the pipeline becomes:
\begin{lstlisting}
DecisionSpec -> PMHD -> ArtifactIR -> Oracle -> String
  -> Write to disk
  -> cargo check        (does it compile?)
  -> cargo test         (do tests pass?)
  -> cargo clippy       (is it idiomatic?)
  -> IF FAIL: feed error to Oracle, iterate
  -> IF PASS: Crystal (now proven to compile and test)
\end{lstlisting}

The crystal now means: ``this code compiles, its tests pass, and it
has been formally validated.'' That is a qualitatively different
guarantee.


\section{Axioms}\label{sec:axioms}

\begin{axiom}[Compilation is Validation]\label{ax:compile}
A code artifact that does not compile is not valid, regardless of
its formal crystal properties. The Foundry adds \texttt{cargo check}
as a \emph{ninth gate} for Rust artifacts. An artifact that passes
the 8 formal gates but fails compilation is rejected.
\end{axiom}

\begin{axiom}[Tests are Proof]\label{ax:tests}
Generated code \textbf{MUST} include tests. A code artifact without
tests is structurally incomplete. The Foundry requires at least one
test per public function/method and verifies that all generated tests
pass.
\end{axiom}

\begin{axiom}[Errors are Feedback]\label{ax:feedback}
Compiler errors and test failures are not terminal — they are
information. The Foundry feeds error output back to the Oracle as
structured correction prompts. Each retry cycle narrows the error
space.
\end{axiom}

\begin{axiom}[Workspace is Context]\label{ax:workspace}
The Forge does not operate in a vacuum. It operates in the context
of a real project directory with existing code, existing types,
existing tests. The Foundry provides this context to the Oracle
so that generated code fits into what already exists.
\end{axiom}


%%% =====================================================================
\part{The Build-Test-Fix Loop}\label{part:loop}

\section{Loop Architecture}\label{sec:looparch}

\begin{algorithm}[H]
\caption{\textsc{Foundry\_Fabricate}: Closed-loop fabrication}
\label{alg:foundry}
\begin{algorithmic}[1]
\REQUIRE ForgeResult $F$ (from C23), ProjectDir $D$, Config $\theta$
\ENSURE FabricationResult (success with working project, or failure with diagnostics)
\STATE $\mathrm{attempt} \gets 0$
\STATE $\mathrm{errors} \gets []$
\REPEAT
  \STATE $\mathrm{attempt} \gets \mathrm{attempt} + 1$
  \STATE \textit{// Step 1: Write files to project directory}
  \STATE $\textsc{WriteProject}(F, D)$
  \STATE \textit{// Step 2: Compile}
  \STATE $\mathrm{compile} \gets \textsc{CargoCheck}(D)$
  \IF{$\mathrm{compile.failed}$}
    \STATE $\mathrm{errors} \gets \mathrm{compile.stderr}$
    \STATE $F \gets \textsc{OracleCorrect}(F, \mathrm{errors}, \text{``compile''})$
    \STATE \textbf{continue}
  \ENDIF
  \STATE \textit{// Step 3: Lint}
  \STATE $\mathrm{lint} \gets \textsc{CargoClippy}(D)$
  \IF{$\mathrm{lint.warnings} > 0$}
    \STATE $F \gets \textsc{OracleCorrect}(F, \mathrm{lint.output}, \text{``lint''})$
    \STATE \textit{// lint warnings don't block, but attempt correction}
  \ENDIF
  \STATE \textit{// Step 4: Test}
  \STATE $\mathrm{test} \gets \textsc{CargoTest}(D)$
  \IF{$\mathrm{test.failed}$}
    \STATE $\mathrm{errors} \gets \mathrm{test.stderr}$
    \STATE $F \gets \textsc{OracleCorrect}(F, \mathrm{errors}, \text{``test''})$
    \STATE \textbf{continue}
  \ENDIF
  \STATE \textit{// Step 5: All passed}
  \RETURN $\textsc{FabricationSuccess}(D, F, \mathrm{attempt})$
\UNTIL{$\mathrm{attempt} \ge \theta.\mathrm{max\_attempts}$}
\RETURN $\textsc{FabricationFailure}(D, \mathrm{errors}, \mathrm{attempt})$
\end{algorithmic}
\end{algorithm}


\section{Correction Prompts}\label{sec:correction}

\begin{definition}[Correction Prompt]\label{def:correction}
When compilation or tests fail, the Foundry constructs a correction
prompt that includes:
\begin{enumerate}[nosep]
  \item The original code that was generated.
  \item The exact compiler/test error output.
  \item The specific file and line where the error occurred.
  \item The error type classification:
        \begin{itemize}[nosep]
          \item \texttt{syntax}: parse error, missing semicolons, brackets
          \item \texttt{type}: type mismatch, wrong generics, missing From impl
          \item \texttt{import}: missing use statements, wrong module paths
          \item \texttt{trait}: missing trait implementation, wrong bounds
          \item \texttt{lifetime}: borrow checker errors
          \item \texttt{test}: assertion failure, panic in test
          \item \texttt{logic}: test passes but wrong output
        \end{itemize}
  \item How many attempts remain (creates urgency for the Oracle to
        get it right).
\end{enumerate}
\begin{lstlisting}
pub struct CorrectionPrompt {
    pub original_code: String,
    pub error_output: String,
    pub error_file: String,
    pub error_line: Option<usize>,
    pub error_class: ErrorClass,
    pub attempt: usize,
    pub max_attempts: usize,
    pub context: WorkspaceContext, // what else exists in the project
}
\end{lstlisting}
\end{definition}

\begin{definition}[Error-Aware Oracle Call]\label{def:errorcall}
The Oracle (C25) receives the correction prompt as an extended
synthesis request:
\begin{lstlisting}
System: "You previously generated code that failed to compile.
Fix the error. Return ONLY the corrected complete file.
Do not explain. Do not apologize. Just fix it."

User: "File: src/handlers.rs
Error (attempt 2/5):
error[E0277]: the trait bound `User: Serialize` is not satisfied
  --> src/handlers.rs:42:5
  |
42 |     Json(user)
  |     ^^^^ the trait `Serialize` is not implemented for `User`

Original code:
<full file content>

Context:
<workspace summary showing User struct in models.rs without #[derive(Serialize)]>"
\end{lstlisting}
\end{definition}

\begin{remark}
The correction prompt is maximally specific. It does not say ``the
code has errors.'' It says exactly what is wrong, exactly where, and
provides the full context needed to fix it. This dramatically reduces
the Oracle's failure rate on retries. Empirically, most compilation
errors are fixed within 1--2 correction cycles.
\end{remark}


\section{Toolchain Integration}\label{sec:toolchain}

\begin{definition}[Toolchain Executor]\label{def:toolchain}
The Foundry shells out to real toolchain commands:
\begin{lstlisting}
pub struct ToolchainExecutor {
    pub cargo_path: String,         // default: "cargo"
    pub rustfmt_path: String,       // default: "rustfmt"
    pub timeout_seconds: u64,       // default: 120
}

pub struct ToolchainResult {
    pub success: bool,
    pub exit_code: i32,
    pub stdout: String,
    pub stderr: String,
    pub duration_ms: u64,
}

impl ToolchainExecutor {
    pub fn cargo_check(&self, dir: &Path) -> ToolchainResult;
    pub fn cargo_build(&self, dir: &Path) -> ToolchainResult;
    pub fn cargo_test(&self, dir: &Path) -> ToolchainResult;
    pub fn cargo_clippy(&self, dir: &Path) -> ToolchainResult;
    pub fn cargo_fmt_check(&self, dir: &Path) -> ToolchainResult;
    pub fn cargo_fmt(&self, dir: &Path) -> ToolchainResult;
    pub fn cargo_doc(&self, dir: &Path) -> ToolchainResult;
}
\end{lstlisting}
\end{definition}

\begin{requirement}[Toolchain Availability]\label{req:toolchain}
The Foundry \textbf{MUST} check toolchain availability at startup.
If \texttt{cargo} is not found, the Foundry operates in
\emph{dry-run mode}: it writes files and skips compilation. The user
is informed that compilation verification is disabled.
\end{requirement}


%%% =====================================================================
\part{Project Scaffolding}\label{part:scaffold}

\section{Complete Project Generation}\label{sec:projgen}

\begin{definition}[ProjectScaffold]\label{def:scaffold}
When the Foundry generates a new project, it produces a complete,
runnable Cargo project:
\begin{lstlisting}
project-name/
  Cargo.toml              # workspace or single crate
  .gitignore              # Rust-standard gitignore
  README.md               # generated from DecisionSpec
  src/
    main.rs or lib.rs     # entry point
    <module>.rs            # one file per atom
  tests/
    integration.rs         # integration tests
  examples/                # usage examples (if applicable)
  migrations/              # SQL migrations (if DB template)
  .env.example             # environment variables
\end{lstlisting}
\end{definition}

\begin{definition}[Cargo.toml Generation]\label{def:cargotoml}
The Foundry generates \texttt{Cargo.toml} from the template's
dependency list and the Oracle's output:
\begin{lstlisting}
pub struct CargoTomlBuilder;

impl CargoTomlBuilder {
    pub fn build(
        name: &str,
        template: &ArchitectureTemplate,
        atoms: &[AtomArtifact],
    ) -> String {
        // - Package metadata from DecisionSpec
        // - Dependencies from template + Oracle hints
        // - Features from template configuration
        // - Workspace structure if multi-crate
    }
}
\end{lstlisting}
\end{definition}

\begin{definition}[README Generation]\label{def:readme}
A README.md is generated from the DecisionSpec and composition tree:
\begin{lstlisting}
# {project-name}

{spec.description}

## Architecture

{composition tree as markdown}

## Getting Started

```
cargo build
cargo test
cargo run  # (if binary)
```

## Modules

{for each atom: name, purpose, public API summary}

## Generated by ISLS Forge

Crystal ID: {crystal_id}
Template: {template_name}
Generated: {timestamp}
\end{lstlisting}
\end{definition}


%%% =====================================================================
\part{Workspace Intelligence}\label{part:workspace}

\section{Understanding Existing Code}\label{sec:understand}

\begin{definition}[WorkspaceAnalyzer]\label{def:analyzer}
The WorkspaceAnalyzer scans an existing project directory and builds
a model of what exists:
\begin{lstlisting}
pub struct WorkspaceAnalyzer;

impl WorkspaceAnalyzer {
    pub fn analyze(dir: &Path) -> Result<WorkspaceModel>;
}

pub struct WorkspaceModel {
    pub project_name: String,
    pub crate_type: CrateType,         // bin, lib, workspace
    pub modules: Vec<ModuleInfo>,
    pub types: Vec<TypeInfo>,
    pub traits: Vec<TraitInfo>,
    pub functions: Vec<FunctionInfo>,
    pub dependencies: Vec<DepInfo>,
    pub test_count: usize,
    pub loc: usize,                    // lines of code
    pub file_tree: Vec<String>,
}

pub struct ModuleInfo {
    pub name: String,
    pub path: String,
    pub public_items: Vec<String>,     // pub fn, pub struct, pub trait
    pub imports: Vec<String>,
    pub loc: usize,
}

pub struct TypeInfo {
    pub name: String,
    pub kind: TypeKind,                // struct, enum, type alias
    pub fields: Vec<String>,
    pub derives: Vec<String>,
    pub module: String,
}

pub struct TraitInfo {
    pub name: String,
    pub methods: Vec<String>,
    pub implementors: Vec<String>,
    pub module: String,
}

pub struct FunctionInfo {
    pub name: String,
    pub signature: String,
    pub is_public: bool,
    pub has_test: bool,
    pub module: String,
}
\end{lstlisting}
\end{definition}

\begin{definition}[Workspace Context for Oracle]\label{def:wscontext}
When the Oracle synthesizes code for an existing project, it receives
a WorkspaceContext:
\begin{lstlisting}
pub struct WorkspaceContext {
    pub summary: String,               // "Axum REST API, 5 modules, 23 types"
    pub relevant_types: Vec<TypeInfo>,  // types the new code needs
    pub relevant_traits: Vec<TraitInfo>,
    pub relevant_imports: Vec<String>,  // import paths for above
    pub existing_patterns: Vec<String>, // code patterns already used
    pub file_being_modified: Option<String>, // current file content
}
\end{lstlisting}
The Oracle prompt includes the workspace context so it generates code
that uses existing types, follows existing patterns, and imports from
existing modules --- not code that reinvents what already exists.
\end{definition}


\section{Incremental Operations}\label{sec:incremental}

\begin{definition}[Incremental Forge Operations]\label{def:incremental}
The Foundry supports incremental operations on existing projects:
\begin{lstlisting}
pub enum FoundryOperation {
    /// Generate a complete new project from DecisionSpec
    NewProject {
        spec: DecisionSpec,
        output_dir: PathBuf,
    },

    /// Add a new module/component to an existing project
    AddComponent {
        project_dir: PathBuf,
        spec: String,           // "add a user authentication module"
        target_module: Option<String>, // where to add it
    },

    /// Implement a trait for an existing type
    Implement {
        project_dir: PathBuf,
        trait_name: String,
        type_name: String,
        hints: Option<String>,
    },

    /// Add an endpoint to an existing API
    AddEndpoint {
        project_dir: PathBuf,
        method: String,         // GET, POST, PUT, DELETE
        path: String,           // "/users/{id}"
        description: String,
    },

    /// Fix a specific file (rewrite with corrections)
    Fix {
        project_dir: PathBuf,
        file_path: String,
        issue: String,          // "doesn't compile", "test fails", etc.
    },

    /// Generate tests for existing code
    GenerateTests {
        project_dir: PathBuf,
        target: Option<String>, // specific module or all
    },

    /// Refactor a module
    Refactor {
        project_dir: PathBuf,
        target: String,
        instruction: String,
    },

    /// Generate documentation
    Document {
        project_dir: PathBuf,
        target: Option<String>,
    },
}
\end{lstlisting}
\end{definition}


%%% =====================================================================
\part{Quality Gates}\label{part:gates}

\section{The Extended Validation Pipeline}\label{sec:extval}

\begin{definition}[Foundry Validation Pipeline]\label{def:foundryval}
For code artifacts, the validation pipeline extends beyond the 8
formal gates:
\begin{lstlisting}
pub struct FoundryValidation {
    // Standard 8-gate (from C5)
    pub formal: FormalValidationResult,     // 8/8 checks

    // Foundry-specific gates
    pub compiles: bool,                     // cargo check passes
    pub tests_pass: bool,                   // cargo test passes
    pub test_count: usize,                  // number of tests
    pub warnings: usize,                    // cargo clippy warnings
    pub formatted: bool,                    // rustfmt check passes
    pub docs_build: bool,                   // cargo doc succeeds

    // Metrics
    pub loc: usize,                         // lines of code
    pub test_coverage_estimate: f64,        // tests/public_functions
    pub compilation_attempts: usize,        // how many tries needed
    pub oracle_tokens_used: usize,
}
\end{lstlisting}
\end{definition}

\begin{requirement}[Minimum Quality for Crystallization]\label{req:minqual}
A code artifact is crystallized only if:
\begin{enumerate}[nosep]
  \item All 8 formal gates pass.
  \item \texttt{cargo check} succeeds (zero compile errors).
  \item \texttt{cargo test} succeeds (all tests pass).
  \item At least one test exists per public function.
  \item Zero Clippy errors (warnings are tolerated).
\end{enumerate}
Artifacts that compile but have failing tests are \emph{not}
crystallized --- they remain as ``draft'' artifacts in the project
directory.
\end{requirement}


%%% =====================================================================
\part{Fabrication Modes}\label{part:modes}

\section{Interactive Mode}\label{sec:interactive}

\begin{lstlisting}[title={Interactive fabrication session}]
$ isls foundry new --domain rust
ISLS Foundry v1.0.0
> What do you want to build?
REST API for a bookmark manager with tags and search

> Template matched: T01 (REST API Service)
> Decomposing into 9 atoms across 3 molecules...
> Forging atom 1/9: router... [Oracle] ... OK (1.2s)
> Forging atom 2/9: middleware... [Memory] ... OK (0.01s)
> Forging atom 3/9: error_handler... [Oracle] ... OK (0.8s)
> Forging atom 4/9: models... [Oracle] ... OK (1.5s)
> Forging atom 5/9: service... [Oracle] ... OK (2.1s)
> Forging atom 6/9: dto... [Derive] ... OK (0.1s)
> Forging atom 7/9: database... [Oracle] ... OK (1.8s)
> Forging atom 8/9: config... [Static] ... OK (0.01s)
> Forging atom 9/9: main... [Static] ... OK (0.01s)

> Writing project to ./bookmark-api/
> Running cargo check... FAIL (attempt 1/5)
  error[E0599]: no method named `fetch_all` on `SqlitePool`
> Feeding error to Oracle... corrected.
> Running cargo check... OK
> Running cargo test... 12 tests, 12 passed
> Running cargo clippy... 0 errors, 2 warnings
> Running cargo fmt... formatted

> Fabrication complete.
  Project: ./bookmark-api/
  Crystal: a7b3c9...
  Files: 12  |  LOC: 847  |  Tests: 12
  Attempts: 2  |  Tokens: 8,420
  Template: T01 (REST API)
  Autonomy: 33% (3/9 from memory/static)
\end{lstlisting}

\section{Batch Mode}\label{sec:batch}

\begin{lstlisting}
# Generate from a DecisionSpec file
isls foundry new --spec bookmark-api.json --output ./bookmark-api/

# Add to existing project
isls foundry add --project ./bookmark-api/ \
  --spec "add a /search endpoint with full-text search on title and tags"

# Fix existing code
isls foundry fix --project ./bookmark-api/ \
  --file src/database.rs --issue "add pagination to list_bookmarks"

# Generate tests for untested code
isls foundry test --project ./bookmark-api/

# Refactor
isls foundry refactor --project ./bookmark-api/ \
  --target src/service.rs --instruction "extract validation into a separate module"
\end{lstlisting}


\section{Headless Mode (for automation)}\label{sec:headless}

\begin{lstlisting}
# JSON in, JSON out (for CI/CD or other tools)
echo '{"operation":"new","spec":{...},"output":"./out/"}' | \
  isls foundry --json

# Returns:
{
  "success": true,
  "project_dir": "./out/",
  "crystal_id": "a7b3c9...",
  "files": 12,
  "loc": 847,
  "tests": 12,
  "attempts": 2,
  "tokens": 8420
}
\end{lstlisting}


%%% =====================================================================
\part{Rust API}\label{part:api}

\begin{lstlisting}
// isls-foundry/src/lib.rs

pub struct Foundry {
    forge: ForgeEngine,
    oracle: OracleEngine,
    catalog: TemplateCatalog,
    toolchain: ToolchainExecutor,
    analyzer: WorkspaceAnalyzer,
    config: FoundryConfig,
}

pub struct FoundryConfig {
    pub max_attempts: usize,          // default: 5
    pub require_tests: bool,          // default: true
    pub require_clippy_clean: bool,   // default: false (warnings ok)
    pub require_fmt: bool,            // default: true
    pub auto_fmt: bool,               // default: true
    pub generate_readme: bool,        // default: true
    pub generate_gitignore: bool,     // default: true
    pub dry_run: bool,                // skip compilation
}

pub struct FabricationResult {
    pub success: bool,
    pub project_dir: PathBuf,
    pub crystal: Option<SemanticCrystal>,
    pub files: Vec<GeneratedFile>,
    pub validation: FoundryValidation,
    pub attempts: usize,
    pub tokens_used: usize,
    pub duration_ms: u64,
    pub template_used: Option<String>,
    pub autonomy_ratio: f64,
}

pub struct GeneratedFile {
    pub path: String,
    pub content: String,
    pub source: SynthesisSource,      // Oracle, Memory, Static, Derive
    pub loc: usize,
    pub test_count: usize,
}

impl Foundry {
    pub fn new(config: FoundryConfig, forge: ForgeEngine,
               oracle: OracleEngine, catalog: TemplateCatalog) -> Self;

    /// Full fabrication: spec -> working project on disk
    pub fn fabricate(
        &mut self, op: FoundryOperation
    ) -> Result<FabricationResult>;

    /// Analyze an existing workspace
    pub fn analyze(&self, dir: &Path) -> Result<WorkspaceModel>;

    /// Check if toolchain is available
    pub fn toolchain_available(&self) -> bool;
}
\end{lstlisting}


%%% =====================================================================
\part{Gateway Integration}\label{part:gateway}

\begin{lstlisting}[title={New Gateway Endpoints}]
POST /foundry/new         Body: {spec, domain, output_dir}
                          -> FabricationResult

POST /foundry/add         Body: {project_dir, spec}
                          -> FabricationResult (incremental)

POST /foundry/fix         Body: {project_dir, file, issue}
                          -> FabricationResult (fix)

POST /foundry/test        Body: {project_dir, target?}
                          -> FabricationResult (test generation)

GET  /foundry/analyze     Params: ?dir=path
                          -> WorkspaceModel

GET  /foundry/status      -> current fabrication progress
\end{lstlisting}


%%% =====================================================================
\part{Acceptance Tests}\label{part:tests}

\begin{longtable}{l l p{5.5cm}}
\toprule
\textbf{ID} & \textbf{Name} & \textbf{Procedure} \\
\midrule
\endfirsthead \midrule \endhead \bottomrule \endfoot
AT-FD1 & New project compiles & Fabricate REST API from spec; verify
  \texttt{cargo check} passes. \\
AT-FD2 & Tests pass & Fabricate project; verify \texttt{cargo test}
  passes with $\ge 1$ test per public fn. \\
AT-FD3 & Compile-fix loop & Mock Oracle to return code with missing
  import on first call, fixed on second; verify 2 attempts,
  success. \\
AT-FD4 & Budget exhaustion & Set max\_attempts=1; mock Oracle to
  return broken code; verify failure with diagnostics. \\
AT-FD5 & Workspace analysis & Create a Rust project with 3 modules;
  run analyzer; verify correct module/type/function counts. \\
AT-FD6 & Incremental add & Create project; add endpoint via Foundry;
  verify new route in router and \texttt{cargo check} still passes. \\
AT-FD7 & Fix operation & Create project with compile error; run
  fix; verify error resolved. \\
AT-FD8 & Test generation & Create project with untested functions;
  run test generation; verify tests created and passing. \\
AT-FD9 & Scaffold completeness & Fabricate project; verify Cargo.toml,
  README.md, .gitignore, src/, tests/ all present. \\
AT-FD10 & Dry run mode & Fabricate with dry\_run=true; verify files
  written but no cargo commands executed. \\
AT-FD11 & Workspace context & Add endpoint to existing project;
  verify Oracle prompt includes existing types and imports. \\
AT-FD12 & Crystal contains validation & Fabricate project; inspect
  crystal; verify it records compilation\_attempts and
  test\_count. \\
AT-FD13 & Formatted output & Fabricate project; verify all .rs files
  pass \texttt{rustfmt --check}. \\
AT-FD14 & README generated & Fabricate project; verify README.md
  contains project description, architecture, and getting started. \\
\end{longtable}


%%% =====================================================================
\part{Handoff}\label{part:handoff}

\section{Test Count}

\begin{tabular}{l r r}
\toprule
\textbf{Phase} & \textbf{New Tests} & \textbf{Cumulative} \\
\midrule
Phase 1--7 + Genesis & 265 & 265 \\
Phase 8 (C27) & 14 & $\ge$ 279 \\
\bottomrule
\end{tabular}

\section{Completion Criteria}

\begin{enumerate}[nosep]
  \item Fabricating a new project from DecisionSpec produces a
        directory with Cargo.toml, src/, tests/, README.md.
  \item Generated code compiles (\texttt{cargo check} passes).
  \item Generated tests pass (\texttt{cargo test} passes).
  \item Compile-fix loop works: Oracle corrects errors on retry.
  \item Workspace analyzer correctly models existing projects.
  \item Incremental operations (add, fix, test, refactor) work
        on existing projects.
  \item Crystal records compilation status and test count.
  \item Dry-run mode works without \texttt{cargo} installed.
  \item Total tests $\ge 279$, zero warnings.
  \item \textbf{27 Crates. The Forge that compiles its own output.}
\end{enumerate}


\vfill
\begin{center}
\rule{0.5\textwidth}{0.4pt}\\[0.5cm]
{\small Generated for \ISLS{} Extension Phase 8 v1.0.0.\\
The Oracle generates. The compiler verifies. The loop closes.\\
Code that doesn't compile is not a crystal. Code that doesn't test\\
is not complete. Code that doesn't fit is not useful.\\[0.3cm]
The Foundry makes all three guarantees.}
\end{center}

\end{document}
